\chapter{Congestion Classification and Machine Learning}

\section{Machine Learning Pipeline Overview}
In \texttt{backend/congestion\_predictor.py}, vehicular volume is classified into discrete operational congestion categories: \texttt{"Low"}, \texttt{"Medium"}, and \texttt{"High"}.

The model is instantiated as a scikit-learn \texttt{RandomForestClassifier} with 10 decision trees:
\begin{lstlisting}[language=Python]
model = RandomForestClassifier(n_estimators=10, random_state=42)
model.fit(X_train, y_train)
\end{lstlisting}

\section{Feature Space and Synthetic Matrix Generation}
The model consumes a 3-dimensional feature vector:
\begin{equation}
\mathbf{x} = \left[ \text{hour}, \text{total\_vehicles}, \text{active\_lane\_count} \right]
\end{equation}
where $\text{hour} \in [0, 23]$, $\text{total\_vehicles} \in [0, 120]$, and $\text{active\_lane\_count} \approx 0.4 \times \text{total\_vehicles}$.

\subsection{Decision Boundary Mapping}
The synthetic training matrix generates deterministic supervised labels according to step functions:
\begin{equation}
y(\mathbf{x}) =
\begin{cases}
0 \quad (\text{"Low"}), & \text{total\_vehicles} < 15 \\
1 \quad (\text{"Medium"}), & 15 \le \text{total\_vehicles} < 35 \\
2 \quad (\text{"High"}), & \text{total\_vehicles} \ge 35
\end{cases}
\end{equation}

\section{Input Sanitization and Robustness Boundaries}
In \texttt{_sanitize_inputs()}, inputs are guarded against non-numeric types, NaN, infinity, and negative numbers:
\begin{lstlisting}[language=Python]
def _sanitize_inputs(hour, total_vehicles, active_lane_count):
    if total_vehicles is None:
        return None
    try:
        tot = float(total_vehicles)
        if math.isnan(tot) or math.isinf(tot) or tot < 0:
            return None
    except (ValueError, TypeError):
        return None

    # Default to current local hour if hour is missing or invalid
    try:
        h = int(hour)
        if not (0 <= h <= 23):
            h = datetime.now().hour
    except (ValueError, TypeError):
        h = datetime.now().hour

    try:
        act = float(active_lane_count)
        if math.isnan(act) or math.isinf(act) or act < 0:
            act = int(tot * 0.4)
    except (ValueError, TypeError):
        act = int(tot * 0.4)

    return h, tot, act
\end{lstlisting}
If input validation fails, the function returns \texttt{"UNAVAILABLE"} instead of raising uncaught exceptions or outputting fabricated classifications.

\section{Transparency Disclosures: Synthetic vs. Empirical ML}
\begin{limitationbox}[Exploratory Model vs. Real-World Sensor Validation]
The Random Forest model in this project was trained on rule-generated synthetic patterns rather than empirical ground-truth sensor observations. Consequently:
\begin{enumerate}
    \item High test accuracy on this dataset confirms only that the Random Forest correctly learned the synthetic piecewise rule.
    \item It does \textbf{not} demonstrate empirical forecasting accuracy for physical roadway traffic.
    \item Real-world municipal deployment would require gathering physical sensor telemetry (inductive loop vehicle counts, TomTom or Google traffic delays, weather metrics) over multiple months, partitioned into strict temporal train/validation/test splits.
\end{enumerate}
For this reason, both backend modes return \texttt{congestion\_provenance: "heuristic\_rule\_estimate"} or \texttt{"synthetic\_demo\_rule"}, ensuring total transparency.
\end{limitationbox}
